\documentclass[a4paper, serif, 10pt]{moderncv}

%% moderncv
% color and style
\moderncvcolor{black}
\moderncvstyle{banking}

%% page layout/margins
\usepackage[top=2.5cm, bottom=2.5cm, left=1.5cm, right=1.5cm]{geometry}

\nopagenumbers{}

%% Babel config for Brazilian Portuguese language
% ref:
% - https://latex3.github.io/babel/guides/locale-portuguese.html
\usepackage[brazilian]{babel}

%% fontspec
\usepackage{fontspec}

\IfFontExistsTF{Linux Libertine}{
  \setmainfont[Ligatures={TeX, Common}, Numbers=OldStyle]{Linux Libertine}
}{}
\IfFontExistsTF{Linux Libertine O}{
  \setmainfont[Ligatures={TeX, Common}, Numbers=OldStyle]{Linux Libertine O}
}{}

%% CTeX
% CJK support, used to show CJK characters in the resume
%
% - heading=false: disable ctex heading
% - punct=kaiming: use kaiming punctuations styles for CJK
% - scheme=plain: use plain scheme, do not override \normalsize font size
% - space=auto: space settings for CJK characters
%
% CJK fonts are resolved per locale language, e.g. Noto Serif/Sans CJK JP
% for Japanese resumes, with fallback to the next region if not installed
%
% ref:
% - http://ctan.mirrorcatalogs.com/language/chinese/ctex/ctex.pdf
\usepackage[UTF8, heading=false, punct=kaiming, scheme=plain, space=auto]{ctex}

\IfFontExistsTF{Noto Serif CJK SC}{
  \setCJKmainfont{Noto Serif CJK SC}
}{}
\IfFontExistsTF{Noto Sans CJK SC}{
  \setCJKsansfont{Noto Sans CJK SC}
}{}

%% line spacing
\usepackage{setspace}
\setstretch{1.125}

%% URL styling - use normal text instead of monospace
\urlstyle{same}

%% Patch \httplink and \httpslink to handle full URLs
% The original moderncv macros always prepend http:// or https:// to URLs.
% This causes issues when the URL already has a protocol (e.g., https://example.com)
% resulting in malformed URLs like https://https://example.com.
% This patch checks if the URL already contains "://" and uses it directly if so.
% Note: We use \str_set:Nx (x-expansion) to expand macros like \@homepage first.
\ExplSyntaxOn
\str_new:N \g_yamlresume_url_str

\RenewDocumentCommand{\httpslink}{O{}m}{
  \str_set:Nx \g_yamlresume_url_str {#2}
  \str_if_in:NnTF \g_yamlresume_url_str {://}
    { \url{#2} }
    { \url{https://#2} }
}

\RenewDocumentCommand{\httplink}{O{}m}{
  \str_set:Nx \g_yamlresume_url_str {#2}
  \str_if_in:NnTF \g_yamlresume_url_str {://}
    { \url{#2} }
    { \url{http://#2} }
}
\ExplSyntaxOff

\name{Mariana Costa Almeida}{}
\title{Engenheira de IA Sênior | LLMs, MLOps e Visão Computacional}
\phone[mobile]{+55 11 98765-4321}
\email{mariana.almeida.ai@example.com}
\homepage{https://marianacosta-ai.dev}

\address{Rua dos Pinheiros, 1234, Apto 56, São Paulo, São Paulo, Brasil, 05422-001}{}{}

\extrainfo{{\small \faLinkedin}\ \href{https://www.linkedin.com/in/marianacostaai}{@marianacostaai} {} {} {} • {} {} {} 
{\small \faGithub}\ \href{https://github.com/marianacostaai}{@marianacostaai} {} {} {} • {} {} {} 
{\small \faMedium}\ \href{https://medium.com/@marianacostaai}{@marianacostaai}}

\begin{document}

\maketitle

\section{Informações Básicas}

\cvline{}{Engenheira de IA com mais de 7 anos de experiência no desenvolvimento e implantação de soluções de aprendizado de máquina em produção. Especializada em modelos de linguagem, visão computacional e pipelines de MLOps, com forte atuação em arquiteturas escaláveis e liderança técnica.
%
\begin{itemize}
\item Projetou e implantou sistemas de recomendação e busca semântica com LLMs, reduzindo o tempo de resposta em 40\%.
\item Experiência sólida em Python, PyTorch, TensorFlow, AWS e Kubernetes, com foco em reprodutibilidade e monitoramento contínuo.
\item Liderou equipes multidisciplinares e programas de mentoria, promovendo boas práticas de IA responsável e experimentação ágil.
\end{itemize}}
\section{Formação}

\cventry{mar. de 2018–dez. de 2020}
        {Mestrado, Ciência da Computação com ênfase em Inteligência Artificial, Nota: 3.9/4.0}
        {Universidade de São Paulo}
        {\url{https://www.usp.br}}
        {}
        {\begin{itemize}
\item Dissertação sobre modelos de linguagem aplicados à sumarização de documentos jurídicos.
\item Publicação em simpósio nacional e colaboração com laboratório de IA da universidade.
\item Monitoria em disciplina de aprendizado profundo por dois semestres.
\end{itemize}
\textbf{Cursos}: Aprendizado de Máquina Avançado, Processamento de Linguagem Natural, Visão Computacional, Aprendizado por Reforço, Otimização Convexa}

\cventry{mar. de 2014–dez. de 2017}
        {Bacharelado, Engenharia de Computação, Nota: 3.7/4.0}
        {Universidade Estadual de Campinas}
        {\url{https://www.unicamp.br}}
        {}
        {\begin{itemize}
\item Formação com foco em sistemas distribuídos, computação em nuvem e fundamentos de IA.
\item Trabalho de conclusão de curso sobre detecção de anomalias em séries temporais.
\item Participação em grupo de estudos de aprendizado de máquina e maratonas de programação.
\end{itemize}
\textbf{Cursos}: Algoritmos e Estruturas de Dados, Sistemas Operacionais, Banco de Dados, Redes de Computadores, Estatística e Probabilidade, Inteligência Artificial}
\section{Experiência Profissional}

\cventry{jan. de 2022–Atual}
        {Engenheira de IA Sênior}
        {Lumina AI Brasil}
        {\url{https://www.luminaai.com.br}}
        {}
        {\begin{itemize}
\item Liderou a construção de uma plataforma de busca semântica e geração aumentada por recuperação (RAG) para o setor jurídico, atendendo mais de 50 mil usuários.
\item Desenvolveu pipelines de ingestão, indexação e avaliação de LLMs com LangChain, FAISS e métricas de qualidade automatizadas.
\item Implantou monitoramento de modelos em produção com Prometheus, Grafana e Evidently, reduzindo incidentes em 35\%.
\item Coordenou equipe de 6 engenheiros e cientistas de dados, definindo arquitetura, padrões de código e roadmap técnico.
\item Promoveu práticas de IA responsável, incluindo revisão de vieses e documentação de modelos.
\end{itemize}
\textbf{Palavras-chave}: LLMs, RAG, MLOps, Liderança, Arquitetura}

\cventry{fev. de 2020–dez. de 2021}
        {Engenheira de Machine Learning}
        {DataSemente Analytics}
        {\url{https://www.datasemente.com.br}}
        {}
        {\begin{itemize}
\item Desenvolveu modelos de propensão de compra e churn para clientes de varejo, com ganho de 18\% em precisão sobre a linha de base.
\item Criou features e pipelines de treinamento em Spark e scikit-learn, processando mais de 2 TB de dados diariamente.
\item Automatizou o deploy de modelos com Docker, Airflow e AWS SageMaker, reduzindo o tempo de entrega de 2 semanas para 3 dias.
\item Colaborou com times de produto e negócio para traduzir requisitos em soluções de IA mensuráveis.
\end{itemize}
\textbf{Palavras-chave}: Machine Learning, Spark, SageMaker, Churn, Varejo}

\cventry{jan. de 2018–jan. de 2020}
        {Cientista de Dados Júnior}
        {Inova Analytics}
        {\url{https://www.inovaanalytics.com.br}}
        {}
        {\begin{itemize}
\item Construiu dashboards e modelos estatísticos para apoiar decisões de marketing e operações.
\item Realizou análises exploratórias e testes A/B em campanhas digitais, identificando oportunidades de otimização.
\item Desenvolveu scripts em Python e SQL para automatizar relatórios recorrentes e reduzir trabalho manual.
\item Participou da migração de dados para ambiente cloud, garantindo qualidade e governança.
\end{itemize}
\textbf{Palavras-chave}: Python, SQL, A/B Testing, Estatística, Cloud}
\section{Idiomas}

\cvline{Português}{Nativo ou Bilíngue \hfill \textbf{Palavras-chave}: Nativo}
\cvline{Inglês}{Proficiência Profissional Fluente \hfill \textbf{Palavras-chave}: TOEFL iBT 112, Inglês técnico}
\cvline{Espanhol}{Proficiência Profissional Limitada \hfill \textbf{Palavras-chave}: Espanhol intermediário}
\section{Competências}

\cvline{Aprendizado de Máquina}{Especialista \hfill \textbf{Palavras-chave}: Python, PyTorch, TensorFlow, scikit-learn, XGBoost}
\cvline{Processamento de Linguagem Natural}{Avançado \hfill \textbf{Palavras-chave}: Transformers, LLMs, RAG, LangChain, Hugging Face}
\cvline{Visão Computacional}{Avançado \hfill \textbf{Palavras-chave}: OpenCV, YOLO, Detectron2, Segment Anything}
\cvline{MLOps e Cloud}{Avançado \hfill \textbf{Palavras-chave}: AWS SageMaker, Docker, Kubernetes, MLflow, Airflow}
\cvline{Engenharia de Dados}{Intermediário \hfill \textbf{Palavras-chave}: Spark, SQL, Kafka, dbt, Snowflake}
\cvline{Liderança Técnica}{Avançado \hfill \textbf{Palavras-chave}: Mentoria, Arquitetura de IA, Agile, IA responsável}
\section{Prêmios}

\cventry{nov. de 2023}
        {Lumina AI Brasil}
        {Prêmio de Inovação em IA}
        {}
        {}
        {Reconhecimento pelo desenvolvimento da plataforma de busca semântica com LLMs, que gerou economia estimada de R\$ 2 milhões ao ano.}

\cventry{out. de 2021}
        {Simpósio Brasileiro de Inteligência Artificial}
        {Melhor Artigo Científico}
        {}
        {}
        {Premiação pelo artigo sobre aprendizado por reforço aplicado à otimização de recursos em nuvem.}
\section{Certificados}

\cventry{mar. de 2022}
        {AWS}
        {AWS Certified Machine Learning – Specialty}
        {\url{https://aws.amazon.com/certification/certified-machine-learning-specialty/}}
        {}
        {}

\cventry{jun. de 2023}
        {Google Cloud}
        {Google Cloud Professional Machine Learning Engineer}
        {\url{https://cloud.google.com/certification/machine-learning-engineer}}
        {}
        {}

\cventry{set. de 2021}
        {TensorFlow}
        {TensorFlow Developer Certificate}
        {\url{https://www.tensorflow.org/certificate}}
        {}
        {}
\section{Publicações}

\cventry{out. de 2021}
        {Aprendizado por Reforço para Otimização de Recursos em Nuvem}
        {Simpósio Brasileiro de Inteligência Artificial (SBIA)}
        {\url{https://sol.sbc.org.br/index.php/sbia}}
        {}
        {\begin{itemize}
\item Propõe um agente de aprendizado por reforço para alocação dinâmica de recursos computacionais.
\item Demonstra redução de custos operacionais em ambientes de nuvem com manutenção de SLA.
\item Discute desafios de escalabilidade e reprodutibilidade em experimentos de longa duração.
\end{itemize}}
\section{Projetos}

\cventry{jan. de 2024–jun. de 2024}
        {Assistente virtual que combina busca semântica e LLMs para responder dúvidas jurídicas com citações.}
        {Assistente Jurídico com RAG}
        {\url{https://github.com/marianacostaai/assistente-juridico-rag}}
        {}
        {\begin{itemize}
\item Construiu pipeline de ingestão e chunking de documentos jurídicos com LangChain e FAISS.
\item Implementou avaliação automatizada de respostas com métricas de fidelidade e cobertura.
\item Desenvolveu API em FastAPI e interface web para demonstração do produto.
\end{itemize}
\textbf{Palavras-chave}: RAG, LLMs, LangChain, FAISS, FastAPI}

\cventry{mar. de 2022–dez. de 2022}
        {Modelo de visão computacional para triagem de exames de imagem com foco em transparência.}
        {Classificador de Imagens Médicas}
        {\url{https://github.com/marianacostaai/classificador-imagens-medicas}}
        {}
        {\begin{itemize}
\item Treinou e avaliou redes neurais convolucionais para classificação de exames de raios-X.
\item Aplicou técnicas de aumento de dados e explicabilidade com Grad-CAM.
\item Publicou resultados em workshop de IA aplicada à saúde.
\end{itemize}
\textbf{Palavras-chave}: Visão Computacional, PyTorch, Grad-CAM, Saúde, CNN}

\cventry{ago. de 2021–fev. de 2022}
        {Infraestrutura reprodutível para treinamento, versionamento e deploy de modelos de recomendação.}
        {Pipeline de MLOps para Recomendação}
        {\url{https://github.com/marianacostaai/pipeline-mlops-recomendacao}}
        {}
        {\begin{itemize}
\item Orquestrou treinamentos com Airflow e versionamento de dados e modelos com MLflow.
\item Configurou deploy em Kubernetes com monitoramento de deriva de dados.
\item Reduziu o tempo de treinamento e validação de modelos em 30\%.
\end{itemize}
\textbf{Palavras-chave}: MLOps, Airflow, MLflow, Kubernetes, Recomendação}
\section{Interesses}

\cvline{Pesquisa em IA}{LLMs, Aprendizado por Reforço, Ética em IA}
\cvline{Esportes}{Corrida, Ciclismo, Natação}
\cvline{Música}{Violão, Piano}
\section{Voluntariado}

\cventry{mar. de 2022–Atual}
        {Mentora Técnica}
        {PyLadies São Paulo}
        {\url{https://www.pyladies.com}}
        {}
        {\begin{itemize}
\item Mentora mulheres em transição de carreira para tecnologia, com foco em Python e fundamentos de IA.
\item Organiza workshops de introdução a machine learning e boas práticas de projetos.
\item Apoia a revisão de currículos e a preparação para entrevistas técnicas.
\end{itemize}}
    
\end{document}